% from ega4-10-fr.tex
\oldpage[IV-3]{112}
ouvertes (\textbf{I}, 6.1.10), et on peut par suite supposer aussi $X$
irréductible. Alors, comme les anneaux locaux de $X$ sont caténaires
($\mathbf 0$, 16.5.12), l'hypothèse $3^\circ$ de (10.6.1) entraîne que $X$ est
biéquidimensionnel ((5.1.5) et ($\mathbf 0$, 14.3.3))~; il suffit donc d'appliquer
(10.8.4).

On notera que si $S$ est le \emph{spectre d'un corps ou de $\mathbf Z$}, il
vérifie les conditions de (10.8.5).

\begin{corollary}[10.8.6]
\label{IV.10.8.6-fr}
Les hypothèses étant celles de (10.8.5), pour tout $x\in X$, le nombre
$\operatorname{prof}^{*}_{x}(\mathcal F)$ est l'unique entier $n$ ayant la propriété
suivante~: il existe un voisinage ouvert $U$ de $x$ dans $X$ tel que pour tout
point $x'\in U\cap\overline{\{x\}}$, fermé dans $U$, on ait
$\operatorname{prof}(\mathcal F_{x'})=n$. En particulier, pour que
$\operatorname{prof}^{*}_{x}(\mathcal F)\geq m$ (resp.
$\operatorname{prof}^{*}_{x}(\mathcal F)\leq m$), il faut et il suffit qu'il existe
un voisinage ouvert $V$ de $x$ dans $X$ tel que, pour tout
$x'\in V\cap\overline{\{x\}}$, fermé dans $V$, on ait
$\operatorname{prof}(\mathcal F_{x'})\geq m$ (resp.
$\operatorname{prof}(\mathcal F_{x'})\leq m$).
\end{corollary}

En effet, on peut se borner au cas où $x$ n'est pas fermé~; si
$\operatorname{prof}^{*}_{x}(\mathcal F)=n$, l'ensemble des $y\in X$ tels que
$\operatorname{prof}^{*}_{y}(\mathcal F)\leq n$ est fermé en vertu de (10.8.5), donc
contient $\overline{\{x\}}$, et en vertu de la semi-continuité inférieure de
$y\mapsto\operatorname{prof}^{*}_{y}(\mathcal F)$, il existe un voisinage ouvert $U$
de $x$ tel que $\operatorname{prof}^{*}_{y}(\mathcal F)=n$ pour tout
$y\in U\cap\overline{\{x\}}$, donc $\operatorname{prof}(\mathcal F_{x'})=n$ si
$x'\in U\cap\overline{\{x\}}$ est fermé dans $U$ (puisque la notion de profondeur
rectifiée est locale).

Pour les préschémas vérifiant les hypothèses de (10.8.5), la notion de profondeur
rectifiée peut donc se définir à l'aide des valeurs de la profondeur aux
\emph{points fermés} de $X$ (ces derniers formant un ensemble \emph{très dense}
dans toute partie fermée de $X$).

\begin{proposition}[10.8.7]
\label{IV.10.8.7-fr}
Soit $S$ un préschéma localement noethérien vérifiant les conditions $1^\circ$,
$2^\circ$ et $3^\circ$ de (10.6.1). Soient $X$ un préschéma localement de type
fini sur $S$, $\mathcal F$ un $\mathcal O_X$-Module cohérent,
$Y=\operatorname{Supp}(\mathcal F)$. Alors, pour tout $x\in Y$, on a
\begin{equation}
\label{IV.10.8.7.1-fr}
\tag{10.8.7.1}
\operatorname{prof}^{*}_{x}(\mathcal F)
=\dim_x(Y)-\operatorname{coprof}(\mathcal F_x).
\end{equation}
\end{proposition}

En effet, par définition, on a
$\operatorname{coprof}(\mathcal F_x)
=\dim(\mathcal F_x)-\operatorname{prof}(\mathcal F_x)$, et il résulte de (10.6.4) que
l'on a $\dim_x(Y)=\dim(\overline{\{x\}})+\dim(\mathcal F_x)$~; d'où (10.8.7.1) par
définition de $\operatorname{prof}^{*}_{x}(\mathcal F)$.

\begin{corollary}[10.8.8]
\label{IV.10.8.8-fr}
Les hypothèses sur $f$ et $\mathcal F$ étant celles de (9.9.1), la fonction
$x\mapsto\operatorname{prof}^{*}_{x}(\mathcal F_{f(x)})$ est constructible.
\end{corollary}

Notons que pour tout $s\in S$, la fibre $X_s$ étant un préschéma localement de
type fini sur $k(s)$, est un préschéma de Jacobson (10.4.7)~; comme en outre
$\operatorname{Spec}(k(s))$ vérifie les conditions de (10.6.1), on a
\[
  \operatorname{prof}^{*}_{x}(\mathcal F_{f(x)})
  =\dim_x(\operatorname{Supp}(\mathcal F_{f(x)}))
  -\operatorname{coprof}((\mathcal F_{f(x)})_x)
\]
(10.8.7). Or, si $Z=\operatorname{Supp}(\mathcal F)$, on a
$Z_{f(x)}=\operatorname{Supp}(\mathcal F_{f(x)})$ (\textbf{I}, 9.1.13) et $Z$ est
localement constructible (8.9.1)~; donc les fonctions
$x\mapsto\dim_x(\operatorname{Supp}(\mathcal F_{f(x)}))$ et
$x\mapsto\operatorname{coprof}((\mathcal F_{f(x)})_x)$ sont localement
constructibles ((9.9.1) et (9.9.3)), ce qui prouve la proposition.

\subsection{Spectres maximaux et ultrapréschémas}
\label{subsection:IV.10.9-fr}

Les résultats de ce numéro ne seront pas utilisés par la suite.

\begin{env}[10.9.1]
\label{IV.10.9.1-fr}
Soit $X$ un préschéma de Jacobson, et soit $S(X)$ l'\emph{espace annelé} dont
l'espace sous-jacent est le sous-espace des points fermés de $X$, et le faisceau
d'anneaux le faisceau induit sur ce sous-espace par $\mathcal O_X$, autrement dit
le faisceau d'anneaux $\psi^*(\mathcal O_X)$, en désignant par
$\psi:S(X)\to X$ l'injection canonique. Comme $\psi$ est un quasi-
